\section{Experiments}
\label{sec:exp}

\paragraph{Protocol.} Unless stated otherwise we use ImageNet-1k
\cite{russakovsky2015imagenet} scores from a SimCLRv2 encoder with a linear probe
\cite{chen2020big}, descriptors taken from a supervised ResNet-50 head \cite{he2016deep} so
that $\varphi$ is exogenous to the scoring model, $\alpha=0.10$, $30\%$ of classes held out
with all of their calibration rows deleted, $n_{\mathrm{cal}}=25$, and five random splits.
Intervals are $95\%$ confidence intervals over splits. Every comparison is made at matched
average set size, as described in \cref{sec:setup}.

\paragraph{Metrics.} \Cref{tab:metrics} defines the quantities we report. The headline
statistic is worst-class coverage, since a single badly covered class is what a
class-conditional guarantee is meant to prevent, but it can only be estimated when classes
have enough evaluation examples. We therefore fix a threshold of $30$ evaluation examples per
class in advance: above the threshold we report per-class worst coverage (regime~A) and below
it a coarser prevalence-binned surrogate (regime~B), and every result states which regime it
falls in. This distinction turns out to matter a great deal, and \cref{sec:depth} returns to
it.

\paragraph{Oracle ceiling.} To put the numbers on a scale we report an oracle that observes
the test labels, forms the exact per-class offset from them, and then selects its own
shrinkage. It is unattainable by construction and is not a method; it exists to say how much
of the available room a method takes. An earlier version of this ceiling used the unshrunk
per-class quantiles and was exceeded by PCC, which is how we found that the unshrunk quantity
is not a bound at all, being $\lambda=1$ applied to a perfect offset at a setting of $\lambda$
that is itself harmful. We report the shrunk ceiling throughout and give the unshrunk values
in \cref{app:oracle}.

Percentages of the ceiling should be read as ratios of averages and not as a per-split
property. On the primary ImageNet arm the ceiling itself varies across the five splits, from
$0.115$ to $0.211$ with a standard deviation of $0.040$, and our own improvement varies from
$0.002$ to $0.158$, so the fraction of the ceiling taken ranges from $2\%$ to $75\%$ split by
split. The same is true of the competitor arm, where the three splits give $67\%$, $20\%$ and
$87\%$. The headline figures quoted below are therefore ratios of means, and the split-level
spread is the honest measure of how well we can pin them down with the number of splits
ImageNet's validation set allows.

% ---------------------------------------------------------------- metrics
\begin{table}[t]
\centering
\small
\setlength{\tabcolsep}{3pt}
\begin{tabular}{@{}ll@{}}
\toprule
Metric & Definition \\
\midrule
$\hat c_y$ (class coverage)
  & $\frac{1}{|J_y|}\sum_{i\in J_y}\mathbf{1}\{Y_i\in\mathcal{C}(X_i)\}$ \\[2pt]
Worst-class (regime A)
  & $\min_{y}\hat c_y$ \\[2pt]
Bin-worst (regime B)
  & $\min_{b}\hat c^{(b)}$, prevalence bins $b$ \\[2pt]
$\Delta$ worst-class
  & worst-class, method $-$ marginal, matched size \\[2pt]
Below target
  & $\frac{1}{|\mathcal{Y}|}\sum_y\mathbf{1}\{\hat c_y<1-\alpha\}$ \\[2pt]
SSCV
  & $\max_{b}\bigl|\hat c^{(b)}-(1-\alpha)\bigr|$, set-size strata \\[2pt]
Worst-slab
  & $\min$ mean $\hat c_y$ over prevalence windows \\[2pt]
Average set size
  & $\frac{1}{N}\sum_i|\mathcal{C}(X_i)|$ \\[2pt]
Empty-set rate
  & $\frac{1}{N}\sum_i\mathbf{1}\{|\mathcal{C}(X_i)|=0\}$ \\
\bottomrule
\end{tabular}
\caption{Evaluation metrics. $J_y$ is the set of test indices with label $y$ and $N$ the
number of test points. Bins and strata are pooled groups, used where a per-class estimate is
too noisy to report: prevalence bins require $200$ test rows each and set-size strata $30$.
The empty-set rate is included because an average set size below one is only possible when a
substantial share of test points receive nothing, a regime in which coverage and set size both
merely look low.}
\label{tab:metrics}
\end{table}

% ---------------------------------------------------------------- competitors
\begin{table}[t]
\centering
\small
\setlength{\tabcolsep}{4pt}
\begin{tabular}{lccc}
\toprule
Method & $\Delta$ worst-class & Set size & Undefined \\
\midrule
\textsc{Standard} \cite{vovk2005algorithmic}   & $+0.000$ & $1.223$ & $0/300$ \\
\textsc{Classwise} \cite{sadinle2019least}     & $+0.000$ & $1.223$ & $\mathbf{300/300}$ \\
\textsc{Clustered} \cite{ding2023class}        & $+0.000$ & $1.223$ & $0/300$ \\
\textsc{Rc3p} \cite{shi2024rank}               & $+0.000$ & $1.223$ & $\mathbf{300/300}$ \\
\textsc{Fuzzy}, quantile \cite{ding2026conformal} & $+0.000$ & $1.223$ & $0/300$ \\
\textsc{Fuzzy}, rarity \cite{ding2026conformal}   & $+0.003$ & $1.223$ & $0/300$ \\
\textsc{Fuzzy}, random \cite{ding2026conformal}   & $+0.027$ & $1.223$ & $0/300$ \\
\midrule
\textbf{PCC (ours)} & $\mathbf{+0.075}$ & $1.223$ & $\mathbf{0/300}$ \\
\midrule
\textit{Oracle ceiling} & \textit{+0.107} & \textit{1.223} & --- \\
\bottomrule
\end{tabular}
\caption{Classes with zero calibration examples: change in worst-class coverage relative to
the marginal threshold, at matched average set size, on the $300$ held-out ImageNet classes,
averaged over three splits. ``Undefined'' counts the classes to which a method could not
assign a threshold, before the pre-registered fallback to the marginal threshold is applied.
Every method is called through its authors' released implementation, and for \textsc{Fuzzy} CP
the kernel bandwidth is swept with the best value kept, which favours the competitor.
Worst-class coverage is decided by the classes a method fails on, so a method can hand most
held-out classes an informative threshold and still score exactly zero.}
\label{tab:competitors}
\end{table}

\subsection{Classes with no calibration examples}
\label{sec:comp}

\Cref{tab:competitors} is our main comparison. PCC improves worst-class coverage on the
held-out classes by $0.075$, which is $70\%$ of the oracle ceiling, while the best competitor
reaches $0.027$ and five of the seven change the statistic by exactly zero. Paired over
splits, PCC exceeds five of the seven with intervals excluding zero; against the two moving
variants of \textsc{Fuzzy} CP the differences are positive, at $0.049$ and $0.072$, but their
intervals include zero at three splits. Three aspects of the table need comment.

\paragraph{Undefined thresholds.} \textsc{Classwise} CP and \textsc{Rc3p} return an infinite
threshold for all $300$ held-out classes, and the \textsc{Rc3p} implementation prints a warning
to that effect for each one. Their entry of zero records the pre-registered fallback to the
marginal threshold, and the methods themselves have nothing to say about these classes, which
is the distinction \cref{tab:methods} is meant to make visible.

\paragraph{Prevalence-based and random embeddings.} \textsc{Fuzzy} CP with the rarity
projection reaches $0.003$ and is beaten by the same method with a random projection, at
$0.027$. Widening thresholds for data-poor classes in an arbitrary direction already buys
something, whereas prevalence adds nothing on top of that. Our own feature ablation reaches
the same conclusion by an independent route, since prevalence-only descriptors drive $\lambda$
to zero in all ten splits (\cref{tab:ablation}). Since random borrowing already helps, the useful
comparison is between one choice of geometry and another, and the descriptors of
\cref{sec:phi} buy $2.8$ times what a random direction buys.

\paragraph{Effect across classes.} The correction lifts the worst class from $0.569$ to
$0.647$ and leaves macro-coverage essentially unchanged, at $0.898$ before and $0.896$ after.
Because average set size is held fixed, that gain is paid for elsewhere: $12\%$ of classes gain
coverage and $27\%$ lose a small amount, with the mean change in each prevalence quintile lying
between $-0.003$ and $+0.001$. Whether the trade is worth making is a question about the
application, which is why we report the minimum and the average side by side.
\Cref{app:whohelped} gives the full breakdown.

% ---------------------------------------------------------------- ablations
\begin{table}[t]
\centering
\small
\setlength{\tabcolsep}{4pt}
\begin{tabular}{lccc}
\toprule
Ablation & $\Delta$ vs.\ full & $d_z$ & Holm \\
\midrule
$\lambda=0$ (no correction)      & $+0.032$ & $1.30$ & \checkmark \\
$\lambda=1$ (unshrunk)           & $+0.584$ & $3.99$ & \checkmark \\
distance descriptors only        & $+0.039$ & $1.49$ & \checkmark \\
prevalence descriptor only       & $+0.032$ & $1.30$ & \checkmark \\
without prevalence               & $+0.011$ & $0.33$ & --- \\
$n^\star$ rule replaced          & $+0.000$ & --- & --- \\
no marginal recalibration        & $+0.000$ & --- & --- \\
\bottomrule
\end{tabular}
\caption{Ablations on ImageNet at $\alpha=0.10$ over ten splits. A positive value means the
full method is better than the ablated one. $d_z$ is Cohen's paired effect size and the last
column marks Holm-corrected significance at level $0.05$ \cite{holm1979simple}. The last two
rows are identical to the full method at matched set size by construction, for the reason
given in the text.}
\label{tab:ablation}
\end{table}

\subsection{Ablations}
\label{sec:ablations}

\Cref{tab:ablation} removes one component at a time over ten splits. Shrinkage carries the
method: both endpoints of the $\lambda$ grid are significantly worse than the selected
interior value, and the unshrunk correction costs $0.584$ in worst-class coverage. Among the
descriptors, using distances alone or prevalence alone is significantly worse than the full
set, while removing prevalence from the full set is not ($p=0.30$), so the geometric features
are what carry the signal.

The final two rows are exactly zero, and the reason is a property of the metric. Marginal recalibration adds one constant to every threshold, and the
matched-size comparison shifts every threshold by one constant, so the two cancel and this
view of the results cannot see that component at all. At the thresholds as they are emitted,
the recalibration step moves marginal coverage from $0.8963$ to $0.8989$ against a target of
$0.90$, at a cost of $0.31$ in average set size, over ten splits. Those raw thresholds are much
wider than the marginal ones, giving an average set size of $2.31$ against $1.21$, which is
precisely what the size matching removes and why the unmatched view cannot be used to compare
methods. We therefore report both views, here and in
\cref{app:invisible}, so that a component the chosen metric cannot see is not recorded as
inert.

% ---------------------------------------------------------------- eval depth
\begin{figure}[t]
\centering
\fbox{\parbox[c][4.2cm][c]{0.92\linewidth}{\centering
\TODO{Figure: evaluation-depth curve. $x$ = evaluation examples per class
(3, 10, 35, 75, all), two panels: (a) $\Delta$ worst-class coverage for PCC with CI band;
(b) the oracle ceiling on the same axis, showing it collapsing to $\le 0$ at 3.
Shade the regime-B region ($<35$) and mark where the four failing settings sit.}}}
\caption{Evaluation depth. Evaluation examples per class are capped inside one dataset with
everything else held fixed. Below roughly $35$ examples per class the oracle ceiling itself
falls to zero, so no method, including one that observes the test labels, can be shown to
improve worst-class coverage. The settings in which our method appeared to fail sit at $35$,
$3$ and $2$ examples per class.}
\label{fig:evaldepth}
\end{figure}

\subsection{The role of evaluation depth}
\label{sec:depth}

Applied without modification, PCC fails on two long-tailed datasets and on two additional
backbones. We traced all four failures to a single axis, and doing so required discarding two
earlier explanations of our own.

\paragraph{Calibration depth.} Capping the number of calibration examples per class inside
ImageNet, with everything else fixed, produces a monotone dose-response: $0.033$ at ten
examples per class rising to $0.060$ at one hundred, with all five intervals excluding zero
(\cref{tab:depth}). Ten examples per class is far thinner than the configurations that fail,
yet the effect is clearly present, so calibration depth scales the effect without explaining
the failures.

\paragraph{Endogenous and exogenous descriptors.} We next hypothesised that the backbone
failures came from using each model's own head as $\varphi$, which makes the descriptors
endogenous to the scores they are used to correct. Running both arms refutes this: with a
fixed exogenous head, whose weights correlate $0.001$ element-wise with each model's own head,
ConvNeXt-T moves from $0.000$ to $0.011$, ViT-B/16 from $-0.029$ to $-0.011$ and ResNet-50
from $0.017$ to $0.006$, with every interval containing zero (\cref{app:backbones}).

\paragraph{Evaluation depth.} \Cref{fig:evaldepth} caps the number of \emph{evaluation}
examples per class, the mirror of the calibration sweep and again inside a single dataset. The
effect is monotone and its interval crosses zero between $35$ and $75$ examples per class.
More importantly, the oracle ceiling collapses along the same axis and is at or below zero by
three examples per class, meaning that a method with access to the test labels could not
improve worst-class coverage there either. In that regime the room to improve has itself
disappeared. The regime threshold of \cref{tab:metrics} also
falls inside this range, since with three examples a class's coverage can take only four
distinct values.

The axis predicts the four failures quantitatively. The additional backbones have $35$
evaluation examples per class and measure $0.017$, which is what the sweep gives at $35$;
Pl@ntNet-300K has three and iNaturalist-2018 has two, where the sweep is at or below zero.

\paragraph{A test of the explanation.} If the explanation is correct then restricting
evaluation to classes that clear the threshold should recover the effect on a long-tailed
dataset. On Pl@ntNet-300K it does. Of the $1081$ classes, $98$ have at least $75$ evaluation
examples, and $34$ of those are held out; over that label space the improvement is $0.019$ with
interval $[0.010,0.028]$ and five of five splits improving. The weaker threshold of $35$
examples keeps more classes but does not reach significance, at $0.020$ with interval
$[-0.011,0.052]$, which is the behaviour \cref{fig:evaldepth} predicts for a slice sitting at
the edge of measurability (\cref{tab:longtail}). The apparent failure on this dataset was an
artefact of what its test set can measure, and the price of removing that artefact is that the
claim now covers under a tenth of the label space.

On iNaturalist-2018 the effect remains exactly zero after the same restriction, and the
reason is not scarcity, since the surviving classes retain $40$ calibration and $56$ evaluation
examples each at the weaker threshold. Instead $\lambda$ is selected as $0$ in all five splits,
so the method returns the marginal threshold unchanged. The shrinkage curve says why. It is not
flat, which would mean $\hat\delta$ is degenerate and every $\lambda$ scores alike; it decreases
monotonically, from $0.695$ at $\lambda=0$ to $0.335$ at $\lambda=1$, so $\hat\delta$ carries
signal and that signal points the wrong way. The regression is correspondingly poor: the
cross-validated error of $g_\theta$ is $0.127$ here against $0.036$ on ImageNet and $0.072$ on
Pl@ntNet-300K. Eight times as many classes fitted in a descriptor space of the same dimension is
the obvious explanation, and richer descriptors are the obvious remedy, but we have not tested
either and offer both as hypotheses. What the shrinkage step does deliver here is that a
$3.5\times$ worse regression costs nothing: the method detects that it cannot help and abstains.

% ---------------------------------------------------------------- held-out
\begin{figure}[t]
\centering
\fbox{\parbox[c][3.6cm][c]{0.92\linewidth}{\centering
\TODO{Figure: held-out fraction curve. $x$ = fraction of classes with zero calibration data
(10--90\%), $y$ = $\Delta$ worst-class coverage with CI band; second $y$-axis or annotation
for \% of oracle ceiling. Mark the $p+2$ floor.}}}
\caption{Amount of supervision required. Worst-class improvement as a function of how many
classes have no calibration data at all. With $100$ labelled classes out of $1000$ the effect
is still $0.052$ with an interval excluding zero. The structural floor is $p+2$ labelled
classes, which is eleven in this configuration.}
\label{fig:heldout}
\end{figure}

\subsection{Amount of supervision required}
\label{sec:supervision}

\Cref{fig:heldout} sweeps the held-out fraction from $10\%$ to $90\%$. The effect peaks at
$50\%$, with $0.080$ or $62\%$ of the ceiling, and remains significant at $90\%$, where only
$100$ classes retain labels, with $0.052$ and interval $[0.007,0.098]$. Four of the five
points have intervals excluding zero; the exception is $10\%$, where only $100$ classes are
held out, so the worst-class statistic computed over them is noisier and the room to improve
is smaller. The floor is structural, as $g_\theta$ is a ridge regression over $p$ features and
needs at least $p+2$ labelled classes, which is eleven in our configuration.

\subsection{Score functions and miscoverage level}
\label{sec:scores}

\Cref{tab:depth} sweeps the score function and $\alpha$. The correction survives all four
scores but is strongest with LAC and attenuates for the adaptive ones. A plausible mechanism
is that APS, RAPS and SAPS already reshape the score using rank or cumulative mass, so part of
the class-difficulty signal that $\delta_y$ captures has been absorbed into the score itself.
At $\alpha=0.01$ the effect is exactly zero, which is consistent with the per-class offsets
becoming indistinguishable from noise as the quantile moves into the extreme tail.

\begin{table}[t]
\centering
\small
\setlength{\tabcolsep}{3.5pt}
\begin{tabular}{llcc}
\toprule
Axis & Setting & $\Delta$ worst-class & \% ceiling \\
\midrule
\multirow{4}{*}{Score}
 & LAC \cite{sadinle2019least}                  & $\mathbf{+0.058}\ [+0.006,+0.110]$ & $42$ \\
 & SAPS \cite{huang2024conformal}               & $+0.031\ [-0.019,+0.081]$ & $24$ \\
 & APS \cite{romano2020classification}          & $+0.022\ [+0.004,+0.040]$ & $15$ \\
 & RAPS \cite{angelopoulos2021uncertainty}      & $+0.010\ [-0.009,+0.029]$ & $10$ \\
\midrule
\multirow{3}{*}{$\alpha$}
 & $0.10$ & $\mathbf{+0.058}\ [+0.006,+0.110]$ & $42$ \\
 & $0.05$ & $+0.040\ [-0.008,+0.088]$ & $39$ \\
 & $0.01$ & $+0.000\ [\ \ 0.000,\ 0.000]$ & $0$ \\
\midrule
\multirow{5}{*}{\shortstack[l]{Cal.\\depth}}
 & $10$/class  & $+0.033\ [+0.016,+0.050]$ & $26$ \\
 & $25$/class  & $+0.030\ [+0.011,+0.049]$ & $23$ \\
 & $50$/class  & $+0.041\ [+0.021,+0.061]$ & $32$ \\
 & $100$/class & $\mathbf{+0.060}\ [+0.032,+0.088]$ & $46$ \\
 & $175$/class & $+0.059\ [+0.031,+0.086]$ & $46$ \\
\bottomrule
\end{tabular}
\caption{Score function, miscoverage level and calibration depth on held-out ImageNet classes
over five splits, and ten splits for calibration depth. The calibration-depth block is the
dose-response referred to in \cref{sec:depth}; the ceiling is flat across it, from $0.128$ to
$0.132$, so the room available is constant across the block and the rise is attributable to
the method.}
\label{tab:depth}
\end{table}

\paragraph{Corrupted inputs.}
\label{sec:shift}
We also calibrate on clean images and evaluate on ImageNet-C \cite{hendrycks2019benchmarking},
over all $15$ corruptions at severities $3$ and $5$ with three splits, for $90$ conditions.
The two slices come from different distributions, so exchangeability fails and the coverage
guarantee is void for every method compared, standard split conformal included. We therefore
read this as a stress test and report the empty-set rate as direct evidence that the guarantee
is gone: at contrast severity $5$, marginal coverage is $0.053$ and $75.7\%$ of test points
receive an empty set. Pooled over conditions PCC gives $0.003$ with interval $[0.001,0.005]$,
so the mechanism survives the shift but the effect is an order of magnitude smaller, and under
severe corruption the whole setting degenerates. Per-corruption numbers are in
\cref{app:imagenetc}.

\paragraph{Computational cost.} PCC trains nothing. The regression is a closed-form ridge fit
over at most $K$ points with $p\le15$ features, and the descriptors are read directly from the
classifier's weights, so no images, gradients or GPU are involved. \Cref{tab:runtime} in the
appendix reports wall-clock times against the competitors, whose kernel-weighted variants
dominate the total cost of our comparison.
